%%%%%%%%%%%%%%%%%%%%%%%%%%%%%%%%%%%%%%%%%%%%%%%%%%%%%%%%%%%%%%%%%%%%%%%%%%%%%%%
% CHAPTER 68: DEFORMATION QUANTIZATION EXAMPLES
%%%%%%%%%%%%%%%%%%%%%%%%%%%%%%%%%%%%%%%%%%%%%%%%%%%%%%%%%%%%%%%%%%%%%%%%%%%%%%%

\chapter{Deformation quantization examples}
\label{chap:deformation-examples}
% Regime III --- Filtered-complete (Convention~\ref{conv:regime-tags}).

Three concrete algebras test the deformation-quantization machinery:
the classical affine Poisson structure on $L\fg^*$, which recovers
$\widehat{\fg}_k$ by quantization
(Example~\ref{ex:classical-affine}); the $\beta\gamma$-system, whose
quartic contact obstruction witnesses the first non-formal quantization
(Corollary~\ref{cor:nms-betagamma-mu-vanishing}); and the Virasoro
coisson algebra, which produces Kontsevich-type weights with
central-charge-dependent corrections.

The organizing principle is the $\Pinf$-chiral algebra
(Definition~\ref{def:pinf-chiral-detailed}): a vertex algebra carrying
\emph{both} a commutative chiral product ($\chirCom$-structure) and a
chiral Lie bracket ($\chirLie$-structure), linked by the Leibniz rule.
Since $\mathsf{Com}^{\mathrm{ch},!} \cong \chirLie$, the two structures sit across
Koszul dual frames, and the $\Pinf$-structure is the algebraic
encoding of a semiclassical limit.  The passage from $\Pinf$ to the
quantized $\Einf$-chiral algebra is then a bar-cobar transport
problem: the quantization map is a Maurer--Cartan element in the
$\Pinf$-deformation complex (the strict model of the corresponding
$L_\infty$-deformation object), and its obstructions are computed by
the chiral Hochschild cohomology $\operatorname{ChirHoch}^*(\cA)$ of
Theorem~\ref{thm:hochschild-bar-cobar}.  The examples below trace this
obstruction theory in each standard family, recording where it
terminates (Heisenberg, affine), where it produces a finite tower
($\beta\gamma$), and where it forces an infinite series of corrections
(Virasoro, $\mathcal{W}_N$).

\section{\texorpdfstring{$\Pinf$ structures}{Pinf structures}}
\label{sec:pinf-structures}
\begin{definition}[\texorpdfstring{$\Pinf$}{P-infinity}-chiral algebra]\label{def:pinf-chiral-detailed}
A \emph{$\Pinf$-chiral algebra} (Definition~\ref{def:pinf-chiral}) carries
\emph{two} compatible chiral-operad structures:
\begin{enumerate}[label=(\roman*)]
\item A $\chirCom$-algebra structure: commutative chiral product (i.e.\
  an $\Einf$-chiral algebra / vertex algebra).
\item A $\chirLie$-algebra structure: chiral Lie bracket $\{-, -\}$.
\item Leibniz compatibility: $\{a, b \cdot c\} = \{a, b\} \cdot c + b \cdot \{a, c\}$.
\end{enumerate}
Since $\mathsf{Com}^{\mathrm{ch},!} \cong \chirLie$, the two chiral structures are
presented across Koszul dual frames.
Equivalently, $\mathcal{A}$ is an algebra over the chiral Poisson operad
$\chirPois = \chirCom \circ \chirLie$.

This differs from a coisson algebra
(Definition~\ref{def:coisson}), which is a commutative
$\mathcal{D}_X$-algebra equipped with a classical $\lambda$-bracket
but no chiral multiplication.  A $\Pinf$-chiral algebra carries two
full chiral operations ($\chirCom$ and $\chirLie$ with Leibniz
compatibility); a coisson algebra carries neither, encoding only
the infinitesimal shadow of the chiral structure.
\end{definition}

\begin{example}[Classical affine Poisson]\label{ex:classical-affine}
Let $\mathfrak{g}^*$ be the dual of a Lie algebra with Kirillov--Kostant Poisson structure:
\[
\{f, g\}(x) = x([df_x, dg_x])
\]
The loop algebra version $L\mathfrak{g}^* = \mathfrak{g}^* \otimes \bC[t, t^{-1}]$ has a
coisson (PVA) structure:
\[
\{J^a_\lambda J^b\} = f^{ab}_c J^c
\]
This is the classical limit ($k \to \infty$) of the Kac--Moody vertex algebra;
its deformation quantization produces $\hat{\mathfrak{g}}_k$ as an $\Einf$-chiral
algebra.
\end{example}


\subsection{Coisson algebras}

\begin{definition}[Coisson algebra]\label{def:coisson}
\index{coisson algebra|textbf}
A \emph{coisson algebra} (Beilinson--Drinfeld \cite{BD04}) is a commutative
$\mathcal{D}_X$-algebra $\mathcal{A}$ equipped with a Lie$^*$ bracket
$\{-,-\}\colon \mathcal{A} \boxtimes \mathcal{A} \to \Delta_* \mathcal{A}$
satisfying the Leibniz rule with respect to the commutative product.
Equivalently, $\mathcal{A}$ is a Poisson vertex algebra (PVA): a commutative
differential algebra with a $\lambda$-bracket $\{a_\lambda b\}$ satisfying
sesquilinearity, skew-symmetry, Jacobi, and Leibniz.

A coisson algebra is \emph{not} a chiral algebra: it carries no OPE, only a
commutative product and a compatible bracket encoding the singular part of a
would-be OPE.  Coisson algebras are the semiclassical limits of vertex algebras
($\Einf$-chiral algebras).
\end{definition}


\begin{theorem}[Coisson = \texorpdfstring{$(\chirPois)^c$}{(chirPois)c}-coalgebra; \ClaimStatusProvedElsewhere{} \cite{BD04}]\label{thm:coisson-coalgebra}
Coisson algebras are exactly coalgebras over the Koszul dual cooperad $(\chirPois)^c$.
\end{theorem}


\section{Quantization formulas}
\label{sec:quantization-formulas}

\subsection{Deformation quantization setup}

\begin{construction}[Quantization map]\label{constr:quantization}
A \emph{deformation quantization} of a $\Pinf$-chiral
algebra $(\mathcal{A}_0, \cdot, \{-,-\})$ is an
$\Eone$-chiral algebra $(\mathcal{A}_\hbar, \star)$ such
that:
\begin{enumerate}[label=(\roman*)]
\item $\mathcal{A}_\hbar = \mathcal{A}_0[[\hbar]]$ as a vector space.
\item $a \star b = a \cdot b + \hbar B_1(a, b) + \hbar^2 B_2(a, b) + \cdots$.
\item $a \star b - b \star a = \hbar \{a, b\} + O(\hbar^2)$.
\end{enumerate}
\end{construction}

\begin{theorem}[Formality for \texorpdfstring{$\Pinf$}{P-infinity}-chiral; \ClaimStatusProvedElsewhere{} \cite{Kon03,FG12}]\label{thm:pinf-formality}
Every $\Pinf$-chiral algebra admits a deformation quantization to an $\Eone$-chiral
algebra. The quantization is unique up to gauge equivalence
(in the $L_\infty$ MC groupoid of $\Convinf$).
\end{theorem}

\begin{proof}
The result follows from the formality of the $E_2$-operad (Kontsevich 1999,
Tamarkin 2003) via the factorization algebra formalism.  Specifically:

\emph{Step~1: Formality.}
The Francis--Gaitsgory formality theorem
(Theorem~\ref{thm:chiral-formality}, \ClaimStatusProvedElsewhere{} \cite{FG12})
establishes an $L_\infty$ quasi-isomorphism between chiral polyvector
fields and chiral Hochschild cochains.  This guarantees that the
deformation problem is unobstructed: the obstruction class in
$H^3_{\chirPois}(\mathcal{A}_0; \mathcal{A}_0)$ vanishes by formality.

\emph{Step~2: Existence.}
By formality, a Poisson structure $\pi \in H^2_{\chirPois}$ lifts to a
Maurer--Cartan element $\alpha = \hbar\pi + O(\hbar^2)$ in the
Hochschild cochain complex.  The resulting star product is:
\[
a \star b = \sum_{n \geq 0} \hbar^n \sum_{\Gamma \in G_{n,2}} w_\Gamma B_\Gamma(a, b)
\]
where $G_{n,2}$ are admissible graphs, $w_\Gamma$ are Kontsevich weights
(integrals over $\FM_n(\mathbb{C})$), and $B_\Gamma$ are bidifferential
operators.  Associativity follows from Stokes' theorem on the boundary
strata of $\FM_n(\mathbb{C})$.

\emph{Step~3: Uniqueness.}
Two quantizations correspond to gauge-equivalent Maurer--Cartan elements
in the deformation complex (Theorem~\ref{thm:mc-quantization}).  The
gauge group acts transitively on the fiber over a fixed first-order
deformation by standard obstruction theory: $H^1$ controls the
infinitesimal gauge transformations, and the formality quasi-isomorphism
ensures no higher obstructions to gauge equivalence.
\end{proof}


\subsection{Explicit star product}

\begin{computation}[Star product through order \texorpdfstring{$\hbar^2$}{2}]\label{comp:star-order-2}
For a $\Pinf$-chiral algebra with Poisson bracket $\{a, b\}$:

\emph{Order $\hbar^0$.} $B_0(a, b) = a \cdot b$.

\emph{Order $\hbar^1$.} $B_1(a, b) = \frac{1}{2}\{a, b\}$.

\emph{Order $\hbar^2$.}  In local coordinates with Poisson bivector $\pi^{ij}$:
\[
B_2(a, b) = \frac{1}{8}\pi^{ij}\pi^{kl}\partial_i\partial_k a \cdot \partial_j\partial_l b + \frac{1}{12}\pi^{ij}(\partial_j\pi^{kl})(\partial_i\partial_k a \cdot \partial_l b - \partial_k a \cdot \partial_i\partial_l b)
\]
The first term corresponds to the ``double Poisson'' Kontsevich graph (both internal vertices connecting directly to the two external vertices); the second involves $\partial\pi$ and vanishes for constant Poisson structures.
\end{computation}


\section{Obstructions and anomalies in examples}
\label{sec:obstructions}

\subsection{Obstruction theory}

\begin{theorem}[Obstruction classes; \ClaimStatusProvedElsewhere{} \cite{Kon03}]\label{thm:obstructions}
The obstruction to quantizing a $\Pinf$-chiral algebra lies in:
\[
\mathrm{Obs} \in H^3_{\chirPois}(\mathcal{A}_0; \mathcal{A}_0)
\]
the third chiral Poisson cohomology.  The obstruction must vanish at
\emph{each order} in $\hbar$ for quantization to exist to all orders; alternatively,
a formality theorem (as in Theorem~\ref{thm:chiral-kontsevich} below) bypasses the
order-by-order obstruction analysis entirely.
\end{theorem}

\begin{example}[No obstruction: affine]\label{ex:no-obs-affine}
For the classical affine Poisson algebra (Example~\ref{ex:classical-affine}):
\[
H^3_{\chirPois}(L\mathfrak{g}^*) = 0
\]
so quantization to $\hat{\mathfrak{g}}_k$ exists for all $k$.
\end{example}

\begin{example}[Obstruction: anomalous theories]\label{ex:anomaly}
Certain physical theories have obstructions:
\begin{enumerate}[label=(\roman*)]
\item Chiral WZW model with ``wrong'' level has $\mathrm{Obs} \neq 0$.
\item The obstruction is the \emph{level quantization condition}: the
      Wess--Zumino class $[H] \in H^3(G, \mathbb{Z})$ requires $k \in \mathbb{Z}$
      for the path integral to be well-defined (arising from $\pi_3(G) \cong \mathbb{Z}$).
\end{enumerate}
\end{example}


\subsection{Anomaly cancellation}

\begin{theorem}[Green--Schwarz mechanism; \ClaimStatusProvedElsewhere{} \cite{Pol98}]\label{thm:green-schwarz}
In string theory, anomalies cancel by introducing a two-form $B$ with modified 
transformation law:
\[
\delta B = \omega_{CS}(\delta A)
\]
where $\omega_{CS}$ is the Chern--Simons form.

At the level of chiral algebras, this modifies the OPE to restore consistency.
\end{theorem}


\subsection{Maurer--Cartan elements and deformations}

\begin{definition}[Maurer--Cartan element]\label{def:mc-element}
For a $\Pinf$-chiral algebra $\mathcal{A}$, a \emph{Maurer--Cartan element} is 
$\alpha \in \mathcal{A}^1$ (degree 1) satisfying:
\[
d\alpha + \frac{1}{2}\{\alpha, \alpha\} = 0
\]
the Maurer--Cartan equation.
\end{definition}

\begin{theorem}[MC elements and quantization; \ClaimStatusProvedElsewhere{} \cite{Kon03,KontsevichSoibelman}]\label{thm:mc-quantization}
Quantizations of $\mathcal{A}_0$ correspond to Maurer--Cartan elements in the 
deformation complex:
\[
\mathrm{Def}(\mathcal{A}_0) = (\mathcal{A}_0[[\hbar]] \otimes \mathfrak{g}_{\chirPois}, d + \hbar\{\alpha, -\})
\]
Two quantizations are gauge-equivalent if and only if their MC elements are related by the
gauge action.
\end{theorem}

\begin{computation}[MC equation in coordinates]\label{comp:mc-coords}
For the classical Heisenberg Poisson algebra with $\{p, q\} = 1$:

The Poisson bivector $\pi = \partial_p \wedge \partial_q$ is the leading-order MC element.  The Kontsevich star product $f \star g = fg + \frac{\hbar}{2}\{f,g\} + O(\hbar^2)$ arises from the formal MC element $\alpha = \hbar\pi + O(\hbar^2)$ in the Hochschild cochain complex.

The MC equation $d\alpha + \frac{1}{2}[\alpha, \alpha] = 0$ at order $\hbar^2$ requires $\frac{1}{2}[\pi, \pi] = 0$, which is the Jacobi identity for the Poisson bracket. Since $\{p,q\} = 1$ is constant, $[\pi,\pi] = 0$ trivially, and $\alpha = \hbar\pi$ is an exact MC element giving the Weyl (Moyal) quantization.
\end{computation}


\section{Explicit deformation quantization computations}
\label{sec:explicit-dq-computations}
\index{deformation quantization!explicit computations}

\subsection{Heisenberg: Moyal product verification}

\begin{computation}[Moyal product as chiral deformation]
\label{comp:moyal-chiral}
\index{Moyal product!chiral}

The Heisenberg algebra $\mathcal{H}_\kappa$ is the
\emph{unique} chiral deformation quantization of the free
Coisson algebra $\mathcal{H}_0 = \bC[a_{-n}]_{n \geq 1}$
with $\lambda$-bracket $\{a_\lambda a\}_0 = \kappa$.

\emph{Classical limit.}
At $\hbar = 0$: the Coisson algebra has commutative product
$a_{-m} \cdot a_{-n} = a_{-m} a_{-n}$ (polynomial multiplication)
and Poisson bracket $\{a_{-m}, a_{-n}\}_0 = m\kappa\delta_{m+n}$
(constant Poisson structure in mode space).

\emph{Quantization.}
The deformed product (the vertex algebra OPE):
\begin{align*}
a_{-m} \star_\hbar a_{-n}
&= a_{-m} a_{-n} + \hbar \cdot m\kappa\delta_{m+n} \cdot 1 \\
&= \normord{a_{-m} a_{-n}} + \hbar \cdot
\text{contraction}
\end{align*}
This is exact at order $\hbar^1$: no higher corrections.
The MC element $\alpha = \hbar \pi$ with
$\pi = \sum_m m\kappa \cdot \partial_{a_{-m}} \wedge
\partial_{a_m}$ is exact (constant Poisson structure).

\emph{Bar complex interpretation.}
The bar complex $\bar{B}(\mathcal{H}_\kappa)$ has
$m_n = 0$ for all $n \geq 3$ (no higher $A_\infty$ operations),
reflecting the exactness of the Moyal product.  The only
nontrivial bar operations are $m_1$ (differential, encoding the modes and conformal grading), $m_2$ (binary coproduct, encoding the OPE $a(z)a(w) \sim \kappa/(z-w)^2$), and $m_0$ (curvature $= \kappa$ at genus~1, the genus-1 deformation).
\end{computation}

\subsection{Kac--Moody as deformation of current algebra}

\begin{computation}[Affine deformation: {\texorpdfstring{$\fg[z] \leadsto
\widehat{\fg}_k$}{[z] _k}}]
\label{comp:affine-deformation}
\index{deformation quantization!affine Kac--Moody}

The affine Kac--Moody algebra $\widehat{\fg}_k$ is a
non-trivial chiral deformation of the current algebra
$\fg[z]$ (the Coisson algebra with $\lambda$-bracket
$\{J^a_\lambda J^b\}_0 = f^{abc}J^c$).

\emph{Order $\hbar^0$:} The commutative product
$\normord{J^a J^b}$ (normal ordering).

\emph{Order $\hbar^1$:}
The first deformation introduces the simple pole
$f^{abc}J^c/(z-w)$ in the OPE.  This is the $\lambda$-bracket
$\{J^a_\lambda J^b\}_1 = f^{abc}J^c$ (already present in the
classical Poisson structure).

\emph{Order $\hbar^2$:}
The central extension $k\delta^{ab}/(z-w)^2$ appears.
This is the genuinely quantum correction: the classical
Poisson bracket $\{J^a, J^b\} = f^{abc}J^c$ admits a central
extension, and the Kontsevich configuration space integral
(Computation~\ref{comp:kontsevich-2vertex}) produces the
coefficient $k$ as a graph weight.

The MC element encoding this deformation:
\[
\alpha = \hbar \sum_{a} J^a \otimes J^a
+ \hbar^2 k \sum_a e^a \otimes e_a
\]
where the first term is the Casimir (encoding the OPE simple pole)
and the second is the central extension (encoding the double pole).

\emph{Associativity obstruction.}
The MC equation at order $\hbar^3$:
\[
d\alpha_3 + \frac{1}{2}[\alpha_1, \alpha_2]
+ \frac{1}{2}[\alpha_2, \alpha_1]
+ \frac{1}{6}[\alpha_1, [\alpha_1, \alpha_1]] = 0
\]
The $[\alpha_1, [\alpha_1, \alpha_1]]$ term involves the Jacobi
identity for $f^{abc}$, which vanishes.  The $[\alpha_1, \alpha_2]$
term involves the action of $f^{abc}$ on the central element $k$,
which vanishes because $k$ is central.  Therefore $\alpha_3 = 0$:
the deformation terminates at order $\hbar^2$.

This reflects the \emph{Koszul} property of $\widehat{\fg}_k$:
the bar complex has $m_n = 0$ for $n \geq 3$ at genus~0.
\end{computation}

\subsection{$\mathcal{W}$-algebras: non-trivial higher-order deformation}

\begin{remark}[Convention: OPE-mode $\lambda$-brackets]
\label{rem:dq-ope-mode-convention}
Throughout this subsection, we write $\lambda$-brackets
in the \emph{OPE-mode convention}:
$\{a_\lambda b\} = \sum_{n \geq 0} \lambda^n \, a_{(n)}b$
(ordinary powers of~$\lambda$).  The rest of the manuscript
uses the \emph{divided-power convention}
$\{a_\lambda b\} = \sum_{n \geq 0} \lambda^{(n)} a_{(n)}b$
with $\lambda^{(n)} = \lambda^n/n!$, in which the coefficient
of~$\lambda^n$ as a polynomial in~$\lambda$ is $a_{(n)}b/n!$.
In particular, the central term $T_{(3)}T = c/2$ appears
as $(c/2)\lambda^3$ here, but as $(c/12)\lambda^3$ in the
divided-power convention.
\end{remark}

\begin{computation}[W-algebra deformation: non-terminating]
\label{comp:w-algebra-deformation}
\index{deformation quantization!W-algebra}
\index{W3 algebra@$\mathcal{W}_3$!deformation}

Unlike the Heisenberg and Kac--Moody examples, the $\mathcal{W}_3$
algebra is a \emph{non-trivial} deformation that does not terminate
at any finite order.

\emph{Classical limit.}
The Coisson algebra $\mathcal{W}_3^{\mathrm{cl}}$ has two
generators $T, W$ with $\lambda$-brackets
(OPE-mode convention; Remark~\ref{rem:dq-ope-mode-convention}):
\begin{align}
\{T_\lambda T\}_0 &= 2T \cdot \lambda + \partial T
\label{eq:w3-poisson-TT} \\
\{T_\lambda W\}_0 &= 3W \cdot \lambda + \partial W
\label{eq:w3-poisson-TW} \\
\{W_\lambda W\}_0 &= \Lambda \cdot \lambda^3
+ \frac{1}{2}\partial\Lambda \cdot \lambda^2
+ \frac{1}{6}\partial^2\Lambda \cdot \lambda
+ (\text{lower order})
\label{eq:w3-poisson-WW}
\end{align}
where $\Lambda = TT - \frac{3}{10}\partial^2 T$ is the
classical composite field.

\emph{Order $\hbar^1$: central charge.}
The first quantum correction introduces the central extension:
\[
\{T_\lambda T\}_1 = \frac{c}{2}\lambda^3
\]
This adds the quartic pole $c/2 \cdot (z-w)^{-4}$ to the $TT$ OPE.

\emph{Order $\hbar^2$: normal ordering correction.}
At order $\hbar^2$, the composite field $\Lambda$ acquires a
quantum correction:
\[
\Lambda_\hbar = \normord{TT}
- \frac{3}{10}\partial^2 T
+ \hbar \cdot \frac{c}{10(5c + 22)}\partial^2 T
\]
The coefficient $c/[10(5c + 22)]$ is the
standard normal-ordering correction to $\Lambda$ itself.
This should not be confused with the OPE structure constant
$\alpha(c) = 16/(22+5c)$ multiplying $\Lambda$ in $W_{(1)}W$
(Computation~\ref{comp:w3-nthproducts}); the two quantities play
different roles.

\emph{Order $\hbar^3$ and beyond.}
At order $\hbar^3$, new corrections appear from triple OPE
contractions in the $W(z)W(w)$ OPE.  These do not vanish
because the Jacobi identity for the $W_3$ algebra involves
the composite field $\Lambda$, which itself depends on $\hbar$.

The MC equation at each order:
\begin{center}
\renewcommand{\arraystretch}{1.2}
\begin{tabular}{c|cl}
Order & Term & Source \\
\hline
$\hbar^1$ & $c/2$ in $TT$ & Central extension \\
$\hbar^2$ & $c/(110+2c)$ in $\Lambda$ & Normal ordering \\
$\hbar^3$ & Jacobi correction & Triple contraction \\
$\hbar^n$ ($n \geq 4$) & Iterated & Higher contractions
\end{tabular}
\end{center}

The non-termination reflects $m_n \neq 0$ for all $n$ in the deformation complex of $\mathcal{W}_3$ ($\mathcal{W}_3$ is Koszul as a chiral algebra by Theorem~\ref{thm:w-algebra-koszul-main}; the non-termination is in the Coisson-to-chiral deformation, not in bar-cobar duality).
\end{computation}

\begin{remark}[Deformation complexity hierarchy]
\label{rem:deformation-complexity}
\index{deformation quantization!complexity hierarchy}
Algebras with linear or constant structure constants ($\mathcal{H}_\kappa$, $\widehat{\fg}_k$, $\beta\gamma$) have terminating deformations ($m_n = 0$ for $n \geq 3$).  Algebras with nonlinear composites ($\mathcal{W}_3$) or non-local $R$-matrices (Yangians) have $m_n \neq 0$ for all $n$.  Koszulness of the bar complex does not constrain deformation complexity.
\end{remark}


\section{Lattice VOA as deformation quantization}
\label{sec:lattice-dq}
\index{lattice VOA!deformation quantization}

\begin{computation}[Lattice cocycle deformation]
\label{comp:lattice-cocycle-deformation}

Let $L$ be an even lattice with bilinear form $\langle \cdot, \cdot \rangle$.
The lattice vertex algebra $V_L$ is constructed from the Heisenberg
algebra $\mathcal{H} = \mathcal{H}(\mathfrak{h})$ (where $\mathfrak{h} = L \otimes_\mathbb{Z} \mathbb{C}$)
and the group algebra $\mathbb{C}[L]$ with a 2-cocycle
$\varepsilon \colon L \times L \to \{\pm 1\}$ satisfying:
\[
\varepsilon(\alpha, \beta) \varepsilon(\beta, \alpha)^{-1}
= (-1)^{\langle \alpha, \beta \rangle}.
\]

At $\hbar = 0$ the algebra is commutative with Poisson bracket from $\langle\cdot,\cdot\rangle$; quantization introduces the vertex operator OPE $e^\alpha(z)e^\beta(w) = \varepsilon(\alpha,\beta)(z-w)^{\langle\alpha,\beta\rangle}e^{\alpha+\beta}(w)+\cdots$, where $\varepsilon$ is determined by the $\overline{\mathrm{FM}}_2(X)$ configuration space integral.  The bar complex $\bar{B}(V_L)$ has degree~1 generators $\mathfrak{h}\oplus L$, with $\varepsilon$ governing the signs of the bar differential.  When $\langle\alpha,\beta\rangle\notin 2\mathbb{Z}$ for some $\alpha,\beta\in L$, the algebra $V_L$ is genuinely $\Eone$-chiral.
\end{computation}

\begin{proposition}[Lattice deformation is one-step; \ClaimStatusProvedHere]
\label{prop:lattice-one-step}
The deformation quantization of the commutative lattice algebra
$V_L^{\mathrm{cl}} = \mathbb{C}[L] \otimes \operatorname{Sym}(\mathfrak{h}[t^{-1}]t^{-1})$
to the vertex algebra $V_L$ terminates at order $\hbar^1$.
Equivalently, the Maurer--Cartan element
$\alpha = \hbar \sum_i e^i \otimes e_i + \hbar \varepsilon$
is exact.
\end{proposition}

\begin{proof}
The bilinear form $\langle \cdot, \cdot \rangle$ is constant on $L$
(it does not depend on the field variables), so the Poisson structure
is constant.  By the same argument as for the Heisenberg algebra
(Computation~\ref{comp:moyal-chiral}), the MC equation $[\alpha, \alpha] = 0$
holds exactly at first order: the Jacobi identity for the constant
Poisson bracket ensures all higher obstructions vanish.

The cocycle $\varepsilon$ contributes signs but no additional deformation
parameters: changing $\varepsilon$ to a cohomologous cocycle $\varepsilon'$
gives an isomorphic vertex algebra $V_L \cong V_L'$ (the isomorphism
class depends only on $[\varepsilon] \in H^2(L, \mathbb{C}^\times)$,
which is determined by $\langle \cdot, \cdot \rangle \bmod 2$).
\end{proof}


\section{Symplectic fermion and logarithmic deformation}
\label{sec:symplectic-fermion-deformation}
\index{symplectic fermion!deformation}

\begin{computation}[Symplectic fermion deformation]
\label{comp:symplectic-fermion-dq}

The symplectic fermion algebra $\mathcal{SF}$ has two generators
$\chi^+, \chi^-$ of conformal weight $1$ with OPE:
\[
\chi^+(z) \chi^-(w) \sim \frac{1}{(z-w)^2}, \qquad
\chi^\pm(z) \chi^\pm(w) \sim 0.
\]
Central charge: $c = -2$.

\emph{Classical limit.}
The Coisson algebra $\mathcal{SF}^{\mathrm{cl}}$ has generators
$\chi^+, \chi^-$ with $\lambda$-bracket
$\{\chi^+_\lambda \chi^-\}_0 = \lambda$ (a \emph{linear}
Poisson bracket, not constant).

\emph{Deformation structure.}
The MC element for the deformation is:
\[
\alpha = \hbar \chi^+ \wedge \chi^-
+ \hbar^2 \cdot 0 + \cdots
\]
The deformation terminates at order $\hbar^1$ because the
Poisson bracket $\{\chi^+, \chi^-\} = \lambda$ is linear
(the structure constants $f^{abc}$ are zero except for
$f^{+-} = 1$), giving $[\alpha, [\alpha, \alpha]] = 0$ by
the Jacobi identity.

The symplectic fermion has logarithmic representations: $L_0$ has a Jordan block on the vacuum module with logarithmic partner $|\omega\rangle = \chi^+_{-1}\chi^-_{-1}|0\rangle$.  Consequently, $\bar{B}(\mathcal{SF})$ has a non-semisimple filtration with Jordan blocks on the $E_1$ page.
\end{computation}

\begin{remark}[Symplectic fermion Koszul dual]
\label{rem:sf-koszul-dual}
The symplectic fermion $\mathcal{SF}$ at $c = -2$ is the $bc$ system at $\lambda = 1$ (Section~\ref{sec:betagamma-koszul-dual}), with Koszul dual $\beta\gamma$ at $\lambda = 1$.  The central charges $c(bc) = -2$ and $c(\beta\gamma) = 2$ follow from the Sugawara construction (or equivalently, from the conformal anomaly of the ghost number current); their sum $c + c' = 0$ is consistent with $bc$--$\beta\gamma$ complementarity.
\end{remark}


\section{Deformation quantization and Koszul duality compatibility}
\label{sec:dq-koszul-compatibility}
\index{deformation quantization!Koszul duality}

\begin{theorem}[Deformation--duality compatibility; \ClaimStatusProvedHere]
\label{thm:dq-koszul-compatible}
Let $\mathcal{A}_0$ be a Coisson algebra and
$\mathcal{A}_\hbar$ its deformation quantization (an $\Einf$-chiral
algebra).  Then:
\begin{enumerate}
\item The Koszul dual of $\mathcal{A}_\hbar$ is a deformation
quantization of the Koszul dual Coisson algebra $\mathcal{A}_0^!$:
\[
(\mathcal{A}_\hbar)^! = (\mathcal{A}_0^!)_{\hbar'}.
\]
\item The deformation parameters are related by the Feigin--Frenkel
involution: if $\mathcal{A}_\hbar = \widehat{\mathfrak{g}}_k$
(affine KM at level $k$), then $\hbar' = -\hbar - 2h^\vee$
(the dual level $k' = -k - 2h^\vee$).
\item The curvature of the bar complex satisfies
the complementarity constraint
(Theorem~\ref{thm:modular-characteristic}\textup{(iv)}):
for affine Kac--Moody and free fields,
$\kappa(\cA) + \kappa(\cA^!) = 0$;
for the Virasoro algebra,
$\kappa(\cA) + \kappa(\cA^!) = 13$; and
for $\mathcal{W}_N$,
$\kappa(\cA) + \kappa(\cA^!) = \varrho_N K_N$
where $\varrho_N = H_N - 1$ and
$K_N = c(\cA) + c(\cA^!)$ is the Koszul conductor.
In each case, the sum is independent of the
deformation parameter~$\hbar$.
\end{enumerate}
\end{theorem}

\begin{proof}
(i) follows from the functoriality of the bar-cobar adjunction
(Theorem~\ref{thm:bar-cobar-isomorphism-main}): the bar construction
$\bar{B}$ is a functor on the category of chiral algebras, and
deformation quantization is a morphism in this category.  The
Koszul dual coalgebra $\bar{B}(\mathcal{A})$, and hence the Koszul
dual algebra $\mathcal{A}^! = \bar{B}(\mathcal{A})^\vee$, therefore
carries the induced deformation.

(ii) For $\widehat{\mathfrak{g}}_k$, the bar complex has curvature
$m_0 = (k + h^\vee)/(2h^\vee) \cdot \kappa_{\mathfrak{g}}$
(Proposition~\ref{prop:km-bar-curvature}).  The Koszul dual
$\widehat{\mathfrak{g}}_{k'}$ has level $k'$ determined by
$(k + h^\vee) + (k' + h^\vee) = 0$, giving $k' = -k - 2h^\vee$.

(iii) The complementarity of curvatures is the content of
Theorem~\ref{thm:quantum-complementarity-main}:
the genus-$g$ obstruction satisfies
$\mathrm{obs}_g(\mathcal{A}) + \mathrm{obs}_g(\mathcal{A}^!)
= H^*(\mathcal{M}_g, \mathcal{Z}(\mathcal{A}))$,
and taking the Euler characteristic gives the $\kappa$-additivity.
\end{proof}

\begin{computation}[Compatibility verification for KM]
\label{comp:dq-koszul-km}
\index{deformation quantization!Kac--Moody compatibility}

For $\widehat{\mathfrak{sl}}_2$ at level $k$:
\begin{itemize}
\item Classical limit ($k \to \infty$):
$\mathcal{A}_0 = \mathfrak{sl}_2[z]$ (current algebra),
$\mathcal{A}_0^! = \mathfrak{sl}_2^*[z]$
(dual current algebra with reversed bracket).
\item Quantization: $\mathcal{A}_k = \widehat{\mathfrak{sl}}_2$ at
level $k$, $\mathcal{A}_k^! = \widehat{\mathfrak{sl}}_2$ at
level $k' = -k-4$.
\item Curvatures: $\kappa(k) = 3(k+2)/4$, $\kappa(k') = 3(-k-4+2)/4 = 3(-k-2)/4$.
Sum: $\kappa(k) + \kappa(k') = 3(k+2)/4 + 3(-k-2)/4 = 0$.
But the curvatures are $\kappa = td/(2h^\vee)$ with $t = k + h^\vee$
and $d = 3$, $h^\vee = 2$, so $\kappa(k) = 3(k+2)/4$ and
$\kappa(k') = 3(k'+2)/4 = 3(-k-2)/4$.

The complementarity gives $c + c' = 6$ (which is $2d = 6$ for
$\mathfrak{sl}_2$), independent of $k$.  $\checkmark$
\end{itemize}
\end{computation}

\begin{computation}[Compatibility verification for Virasoro]
\label{comp:dq-koszul-vir}
For $\mathrm{Vir}_c$ via DS reduction of $\widehat{\mathfrak{sl}}_2$ at level $k$: the Koszul dual is $\mathrm{Vir}_{26-c}$ (Theorem~\ref{thm:w-algebra-koszul-main}), with curvatures $\kappa(c)=c/2$, $\kappa(c')=(26-c)/2$, summing to $13$ independently of $c$.  At $c=0$ the bar complex is uncurved while the dual is maximally curved; at $c=26$ the roles reverse, giving the algebraic uncurved-shadow configuration that appears in the no-ghost discussion.
\end{computation}

\begin{remark}[Filtered deformation quantization]
\label{rem:filtered-dq}
\index{filtered deformation}
The deformation parameter $\hbar$ induces a filtration on the bar complex with $d_{\bar{B}} = d_0 + \hbar\, d_1 + \hbar^2\, d_2 + \cdots$.  The associated spectral sequence has $E_0 =$ classical bar complex, $E_1 = H^*(E_0, d_0)$, and $E_2 = H^*(E_1, d_1)$.  For Koszul algebras the spectral sequence collapses at $E_2$; for non-Koszul algebras higher differentials may be non-trivial.
\end{remark}


\section{Superpotential truncation of \texorpdfstring{$A_\infty$}{A-infinity} operations}
\label{sec:superpotential-ainfty-truncation}
\index{superpotential!A-infinity truncation@$A_\infty$ truncation}
\index{Landau--Ginzburg!A-infinity truncation@$A_\infty$ truncation}

\providecommand{\dCrit}{\mathrm{dCrit}}

\begin{construction}[Superpotential $A_\infty$ structure; \ClaimStatusProvedElsewhere{} \cite{DNP25}]
\label{constr:superpotential-ainfty}
\index{derived critical locus!chiral enhancement}
Let $(X, \psi)$ be a single $n=1$ chiral multiplet with superpotential
$W(X) = X^d/d!$, where $d \geq 2$.  The Koszul dual algebra
$\mathcal{A}^! = \bC[X_n, \psi_n]_{n \geq 0}$ carries a differential
and higher $A_\infty$ operations determined by $W$:
\begin{align}
Q(\psi_n) &= -\frac{1}{(d-1)!}
  \sum_{\substack{m_1 + \cdots + m_{d-1} = n+2-d \\ m_i \geq 0}}
  X_{m_1} \cdots X_{m_{d-1}},
\label{eq:superpotential-Q} \\
m_k(\psi_{n_1}, \ldots, \psi_{n_k}) &= -\frac{1}{k!\,(d-k)!}
  \sum_{\substack{m_1 + \cdots + m_{d-k} = n_1 + \cdots + n_k + k + 1 - d \\ m_i \geq 0}}
  X_{m_1} \cdots X_{m_{d-k}}.
\label{eq:superpotential-mk}
\end{align}
The key structural consequence: $m_k = 0$ for all $k \geq d$.  The degree
of the superpotential controls the $A_\infty$ truncation depth.
\end{construction}

\begin{example}[$d = 2$: massive chiral, $W = \tfrac{1}{2}mX^2$]
\label{ex:superpotential-d2}
\index{massive chiral!bar complex}
Only $m_1 = Q$ is nonzero: the algebra is strictly differential-graded with
no higher operations.  The differential $Q(\psi_n) = -X_{n}$ pairs
generators bijectively.  The bar complex $\bar{B}(\mathcal{A}^!)$ is
acyclic: its cohomology vanishes in all degrees, reflecting the
contractibility of the massive theory.
\end{example}

\begin{example}[$d = 3$: cubic Landau--Ginzburg, $W = \tfrac{1}{6}X^3$]
\label{ex:superpotential-d3}
\index{Landau--Ginzburg!cubic}
The operations $m_1$ and $m_2$ are nonzero; all $m_k = 0$ for $k \geq 3$.
The algebra $\mathcal{A}^!$ is a genuine dg algebra (not merely
$A_\infty$).  This is the XYZ model: its derived critical locus
$\dCrit(W)$ is a classical derived scheme, and $\mathcal{A}^!$ is its
chiral enhancement.
\end{example}

\begin{example}[$d = 4$: quartic potential, $W = \tfrac{1}{24}X^4$]
\label{ex:superpotential-d4}
\index{Landau--Ginzburg!quartic}
Now $m_1$, $m_2$, and $m_3$ are all nonzero: this is the first genuinely
$A_\infty$ case.  The ternary operation $m_3$ encodes the leading-order
departure from classical geometry.  Physically, this corresponds to a
single-field Landau--Ginzburg model with quartic interaction.
\end{example}

\begin{proposition}[Chiral enhancement of the derived critical locus;
\ClaimStatusProvedElsewhere{} \cite{DNP25}]
\label{prop:chiral-dcrit}
\index{derived critical locus!chiral enhancement}
The algebra $\mathcal{A}^!$ of Construction~\ref{constr:superpotential-ainfty}
is the chiral enhancement of the derived critical locus
$\dCrit(W) = \operatorname{Spec}\bigl(\bC[X]/(\partial W)\bigr)^L$.
The $A_\infty$ operations $\{m_k\}_{1 \leq k < d}$ encode the full
homotopy-coherent structure of $\dCrit(W)$ as a chiral algebra; the
truncation at $k = d$ is a direct consequence of $W^{(d+1)} = 0$.
\end{proposition}

\begin{remark}[Bar spectral sequence visibility]
\label{rem:superpotential-bar-ss}
\index{bar complex!superpotential truncation}
The $A_\infty$ truncation by $\deg(W)$ is visible in the bar spectral
sequence of $\bar{B}(\mathcal{A}^!)$: on the $E_1$ page, the generators
contributed by $m_k$ vanish for $k \geq d$, so the spectral sequence has
nonzero differentials only up to page $E_{d-1}$.  For $d = 2$ it
collapses at $E_1$ (acyclicity); for $d = 3$ at $E_2$ (the classical
derived category); and for $d \geq 4$ the pages $E_3, \ldots, E_{d-1}$
carry genuinely $A_\infty$ information invisible to classical algebraic
geometry.  This hierarchy gives a concrete filtration-theoretic meaning
to the slogan ``the degree of $W$ controls the homotopy depth of $\dCrit(W)$.''
\end{remark}

\begin{remark}[Three-pillar interpretation: deformation quantization]
\label{rem:deformation-three-pillar}
\index{deformation quantization!three-pillar interpretation}
In the three-pillar architecture
(\S\ref{sec:concordance-three-pillars}):
(i)~the superpotential truncation ($m_k = 0$ for $k \geq d$)
is a strict $\mathrm{Ch}_\infty$ phenomenon: the bar complex of
$\dCrit(W)$ has secondary Borcherds operations
$j'_{(p_1,\ldots,p_n)} = 0$ for $n \geq d$.  The degree of $W$
controls the homotopy depth of the $\mathrm{Ch}_\infty$-structure;
(ii)~the convolution $sL_\infty$-algebra
$\operatorname{hom}_\alpha(\cC^{\textup{!`}}, \dCrit(W))$ has
transferred brackets $\ell_k^{\mathrm{tr}} = 0$ for $k \geq d$,
and the star product MC element $\Theta_{\star} =
\sum_{n \geq 1} \hbar^n \Theta_n$ is a finite-order projection of
the universal $\Theta_\cA$ via the convolution identification
(Theorem~\ref{thm:convolution-master-identification});
(iii)~the spectral sequence page at which collapse occurs
($E_{d-1}$) corresponds to the codimension filtration on the
log-FM boundary: higher-order collisions contribute at higher
spectral sequence pages, and the truncation at degree~$d$ kills
the codimension-$\geq d$ strata.
\end{remark}
